\section{Arquitectura reproducible}

El proyecto separa configuración, lógica reusable, runners y outputs. Los checkpoints no se
definen por una secuencia informal de notebooks, sino por contratos versionados.

\begin{longtable}{p{0.24\linewidth}p{0.65\linewidth}}
\toprule
Componente & Rol \\
\midrule
\endhead
\artifact{configs/} & contratos YAML, grids, paths, thresholds e invariantes. \\
\artifact{src/geocebada/data/} & loaders y normalización de fuentes. \\
\artifact{src/geocebada/features/} & construcción determinista de capas de features. \\
\artifact{src/geocebada/evaluation/} & split logic, estimadores y métricas de checkpoints. \\
\artifact{src/geocebada/visualization/} & figuras derivadas de artefactos numéricos. \\
\artifact{tools/} & entry points reproducibles para ejecutar cada checkpoint. \\
\artifact{reports/} & métricas, predicciones OOF/pseudo-test, figuras y candidatos. \\
\artifact{docs/} & doctrina, hallazgos canónicos, fuentes e historial. \\
\bottomrule
\end{longtable}

\section{Mapa de ejecución}

La secuencia principal puede reproducirse conceptualmente con:

\begin{lstlisting}[language=bash]
python tools/audit_data_contract_v2.py
python tools/build_parcel_feature_table.py

python tools/build_agronomic_features_v1.py
python tools/run_checkpoint_03b.py
python tools/run_checkpoint_03c.py
python tools/run_checkpoint_03c2.py
python tools/run_checkpoint_03c3.py

python tools/run_checkpoint_04a.py
python tools/run_checkpoint_04b.py
python tools/run_checkpoint_04d1.py
python tools/run_checkpoint_04e1.py
python tools/run_checkpoint_04c1.py
python tools/run_checkpoint_04f.py
python tools/run_checkpoint_05.py

# Additive global-capacity closure, executed after Checkpoint 05:
python tools/run_checkpoint_03d.py
\end{lstlisting}

El orden 04E/04C anterior refleja el orden real de implementación de las hipótesis finales;
la numeración conceptual mantiene 04C como rama global, 04D como local/grafo, 04E como
evidencia externa y 04F como congelamiento.

\section{Artefactos canónicos por checkpoint}

\subsection*{Checkpoint 01}
\begin{itemize}
\item \artifact{configs/data_contract_v2.yaml}
\item \artifact{reports/data_audit_v2.json}
\item \artifact{reports/data_audit_v2.md}
\item \artifact{data/integration_fixture/}
\end{itemize}

\subsection*{Checkpoint 02}
\begin{itemize}
\item \artifact{data/processed/features_v1/parcel_features_all.csv}
\item \artifact{data/processed/features_v1/feature_manifest.json}
\item \artifact{reports/checkpoint_02/cv_folds.csv}
\item \artifact{reports/checkpoint_02/baseline_summary.csv}
\item \artifact{reports/checkpoint_02/checkpoint_02_report.md}
\end{itemize}

\subsection*{Checkpoint 03}
\begin{itemize}
\item \artifact{data/processed/agronomic_features_v1/}
\item \artifact{data/processed/empirical_features_v1/}
\item \artifact{reports/checkpoint_03b/expression_discovery_report.md}
\item \artifact{reports/checkpoint_03c/checkpoint_03c_report.md}
\item \artifact{reports/checkpoint_03c2/checkpoint_03c2_report.md}
\item \artifact{reports/checkpoint_03c3/}
\item \artifact{reports/checkpoint_03d/}
\item \artifact{docs/CHECKPOINT_03D_FINDINGS.md}
\end{itemize}

\subsection*{Checkpoint 04}
\begin{itemize}
\item \artifact{reports/checkpoint_04a/}
\item \artifact{reports/checkpoint_04b/pseudo_competition_splits.csv}
\item \artifact{reports/checkpoint_04d1/}
\item \artifact{reports/checkpoint_04e1/}
\item \artifact{reports/checkpoint_04c1/}
\item \artifact{reports/checkpoint_04f/final_predictions.csv}
\item \artifact{reports/checkpoint_04f/final_prediction_diagnostics.csv}
\end{itemize}

\subsection*{Checkpoint 05}
\begin{itemize}
\item \artifact{reports/checkpoint_05/final_predictions.csv}
\item \artifact{reports/checkpoint_05/final_prediction_diagnostics.csv}
\item \artifact{reports/checkpoint_05/checkpoint_05_report.json}
\item \artifact{docs/CHECKPOINT_05_FINDINGS.md}
\end{itemize}

\subsection*{Checkpoint 03D deployment}
\begin{itemize}
\item \artifact{reports/checkpoint_03d/finalists.csv}
\item \artifact{reports/checkpoint_03d/robustness_summary.csv}
\item \artifact{reports/checkpoint_03d/checkpoint_03d_report.json}
\item local, gitignored: \artifact{models/final/checkpoint03d_global.onnx}
\item local, gitignored: \artifact{models/final/checkpoint03d_global.onnx.json}
\item local, gitignored: \artifact{models/final/checkpoint03d_global.joblib}
\end{itemize}

\section{Invariantes que deben comprobarse antes de una entrega}

\begin{enumerate}
\item El archivo final tiene 59 filas de datos y exactamente dos columnas.
\item \artifact{ID_POLIGONO} es único y coincide exactamente con el set oficial
\artifact{PREDICCION}.
\item No existe ningún target observado en esas 59 filas dentro del input de modelado.
\item La columna final se denomina exactamente \artifact{RENDIMIENTO_T_HA}.
\item Las predicciones Local04D coinciden con los candidatos congelados de 04D.1/04E.1.
\item El CSV final conserva precisión numérica completa.
\item Cualquier versión redondeada del reporte o de la app se trata sólo como presentación.
\end{enumerate}

\section{GPU y determinismo}

CatBoost se ejecuta en GPU cuando corresponde. Las seeds y configuraciones se congelan en YAML.
El predictor final LocalRidge es determinista dados:
\begin{itemize}
\item la misma tabla C4;
\item las mismas coordenadas/IDs;
\item $k=24$, $\alpha=30$, $p=1$;
\item el mismo conjunto de 138 etiquetas.
\end{itemize}

El uso de GPU no forma parte de la semántica del modelo final. Fue una decisión de eficiencia
para los benchmarks CatBoost.

\section{Separación entre artefacto generado y narrativa}

Las tablas de este documento se transcriben de artefactos congelados, pero no constituyen la
fuente de verdad primaria. Si el texto y un CSV/JSON generado difieren, debe investigarse el
runner y corregirse el reporte; nunca debe editarse el output numérico para que coincida con la
narrativa.

\section{Versionado y provenance}

Cada reporte generado registra timestamp y commit cuando el runner lo soporta. La cadena de
provenance final es:
\[
\text{source data}
\rightarrow
\text{Feature Table v1}
\rightarrow
\text{C4 deterministic}
\rightarrow
\text{04B memberships}
\rightarrow
\text{04D LocalRidge}
\rightarrow
\text{04F frozen rule}
\rightarrow
\text{05 promotion gate}
\rightarrow
\text{05 canonical CSV}.
\]

Los Checkpoints~04E y 04C son ramas de validación de hipótesis: pueden generar predicciones
candidatas, pero no alteran silenciosamente la cadena final si sus gates no justifican el
cambio.

\section{Política de outputs pesados}

Los reportes, PDFs y figuras son outputs derivados. En el repositorio, su regeneración pesada
no debe dispararse automáticamente por cada push; los workflows costosos deben ejecutarse de
forma manual y subir artifacts, mientras CI ligero puede revisar tests, lint y estructura.
Esta separación evita convertir una actualización documental en un job de cómputo
innecesario.

\section{Recompilación de este reporte}

Desde \artifact{reporte/tecnico/}:
\begin{lstlisting}[language=bash]
latexmk -pdf -interaction=nonstopmode main.tex
\end{lstlisting}

La bibliografía usa BibLaTeX/Biber. Si no existe latexmk:
\begin{lstlisting}[language=bash]
pdflatex main.tex
biber main
pdflatex main.tex
pdflatex main.tex
\end{lstlisting}

Las figuras usan rutas relativas a \artifact{../../reports/}. La macro
\artifact{\\safefigure} deja un placeholder si una figura no está disponible, permitiendo que
el texto siga compilando en checkouts parciales.


\section{Incidentes reproducibles de los runs finales}

En Checkpoint~05, un nested fold de 18 observaciones hizo inválido
\artifact{PLS(n_components=20)}. La corrección limita componentes al mínimo entre número de
features y tamaño del menor training fold interno. Es una restricción de factibilidad, no
selección basada en target.

En Checkpoint~03D, la primera configuración wide proyectó varios días de cómputo. El run
balanced preservó outer e inner CV y redujo sólo el número de candidatos pesados. Para impedir
mezclar partials incompatibles, \artifact{--resume} verifica un fingerprint del contrato.

El preflight ONNX detectó incompatibilidades reales de conversores. Se fijó
\artifact{onnx<1.22} con \artifact{skl2onnx==1.20.0}; CatBoost usa exportador nativo y
XGBoost/LightGBM negocian el opset máximo soportado. HistGB permanece científicamente válido
aunque no sea deployable con ese stack. El finalista XGBoost exportado pasó el round trip con
diferencia absoluta máxima $3.814697\times10^{-6}$.
